% Generated by `taskgen course examples` from crypto.lfsr (seed 7).
% Do not edit: change examples.yaml and regenerate.
\begin{example}[период и $n$-й член последовательности РСЛОС]\label{ex:lfsr-period}
Последовательность $s_0, s_1, \dots$ над $\mathbb{F}_2$ порождается РСЛОС с характеристическим многочленом \[f(x) = x^{4} + x + 1\] и начальным состоянием $(s_0, \dots, s_{3}) = (1, 0, 0, 1)$. (а)~Выписать рекуррентное соотношение и сопровождающую матрицу $A$ ($S_{n+1} = A S_n$, $S_n = (s_n, \dots, s_{n+3})^T$). (б)~Найти период последовательности. Является ли $f$ неприводимым, примитивным? (в)~Вычислить $s_{96}$ как первую координату вектора $A^{96} S_0$ (бинарный алгоритм возведения матрицы в степень).

\par\noindent\textit{Решение.} (а) $s_{n+4} = s_{n+1} + s_{n}$; \[A = \begin{pmatrix}0 & 1 & 0 & 0 \\ 0 & 0 & 1 & 0 \\ 0 & 0 & 0 & 1 \\ 1 & 1 & 0 & 0\end{pmatrix}.\] (б) Члены последовательности, пока состояние не вернётся к $S_0$: \texttt{10011 01011 11000 1001}, т.\,е. $S_{15} = S_0$, и период $T = 15$. Разложение: $f(x) = x^{4} + x + 1$ неприводим (нет делителей степени $\le 2$). Для неприводимого $f$ период любого ненулевого состояния равен $\operatorname{ord}_f(x)$ и делит $2^{4} - 1 = 15$; $T = 2^{4} - 1$, т.\,е. $x$ порождает $\mathbb{F}_{2^{4}}^*$. Многочлен $f$ примитивный.\par (в) $96 = (1100000)_2$. Справа налево: $v \gets S_0$; при $n_i = 1$ $v \gets A^{2^i} v$; затем $A^{2^{i+1}} = (A^{2^i})^2$:\[\begin{array}{c|c|c|c}i & n_i & A^{2^i} & v \\ \hline 0 & 0 & \left(\begin{smallmatrix}0 & 1 & 0 & 0 \\ 0 & 0 & 1 & 0 \\ 0 & 0 & 0 & 1 \\ 1 & 1 & 0 & 0\end{smallmatrix}\right) & (1, 0, 0, 1)^T \\ \noalign{\smallskip} 1 & 0 & \left(\begin{smallmatrix}0 & 0 & 1 & 0 \\ 0 & 0 & 0 & 1 \\ 1 & 1 & 0 & 0 \\ 0 & 1 & 1 & 0\end{smallmatrix}\right) & (1, 0, 0, 1)^T \\ \noalign{\smallskip} 2 & 0 & \left(\begin{smallmatrix}1 & 1 & 0 & 0 \\ 0 & 1 & 1 & 0 \\ 0 & 0 & 1 & 1 \\ 1 & 1 & 0 & 1\end{smallmatrix}\right) & (1, 0, 0, 1)^T \\ \noalign{\smallskip} 3 & 0 & \left(\begin{smallmatrix}1 & 0 & 1 & 0 \\ 0 & 1 & 0 & 1 \\ 1 & 1 & 1 & 0 \\ 0 & 1 & 1 & 1\end{smallmatrix}\right) & (1, 0, 0, 1)^T \\ \noalign{\smallskip} 4 & 0 & \left(\begin{smallmatrix}0 & 1 & 0 & 0 \\ 0 & 0 & 1 & 0 \\ 0 & 0 & 0 & 1 \\ 1 & 1 & 0 & 0\end{smallmatrix}\right) & (1, 0, 0, 1)^T \\ \noalign{\smallskip} 5 & 1 & \left(\begin{smallmatrix}0 & 0 & 1 & 0 \\ 0 & 0 & 0 & 1 \\ 1 & 1 & 0 & 0 \\ 0 & 1 & 1 & 0\end{smallmatrix}\right) & (0, 1, 1, 0)^T \\ \noalign{\smallskip} 6 & 1 & \left(\begin{smallmatrix}1 & 1 & 0 & 0 \\ 0 & 1 & 1 & 0 \\ 0 & 0 & 1 & 1 \\ 1 & 1 & 0 & 1\end{smallmatrix}\right) & (1, 0, 1, 1)^T \\ \noalign{\smallskip} \end{array}\]$A^{96} S_0 = S_{96} = (1, 0, 1, 1)^T$. Проверка по периоду: $96 \equiv 6 \pmod{15}$, $s_{96} = s_{6}$.\[T = 15,\quad \text{$f$ примитивный},\quad s_{96} = 1.\]
\end{example}
